\begin{abstract}
本文研究六区域异构 AI 任务、可再生能源、储能与电网的协同调度。针对问题一，采用严格的训练（0--2351）、验证（2352--2375）和测试（2376--2399）时间切分，在季节朴素模型与 $3\times3$ Huber 参数网格间逐序列选型，并执行闭环递推；18 条序列的平均 MAE、RMSE 和 WAPE 分别为 25.5892~GPU、34.5070~GPU 和 0.8231。针对问题二，构造成本、碳、新能源、服务与平衡五类增量邻域，各评价 $10000$ 个候选并对最终 $50000$ 个任务执行全量约束审计；平衡方案净购售电成本为 $-421801727.45$~元，碳排放为 $16.5160$~tCO$_2$，新能源利用率为 $0.679962$。针对问题三，在含 no-action 的七个储能策略间用离散 SOC 动态规划搜索，经轨迹去重和支配检验选择平衡方案，其成本为 $-458584632.33$~元、碳排放为 $0.2241$~tCO$_2$。针对问题四，对 17 个情景逐一从共同基线重启，每轮根据上一轮储能轨迹重建 600 个任务候选，最多 8 轮、连续 2 轮无改善停止；共评价 $22200$ 个候选，基准情景实时 SLA 与按期完成率均为 $1$。两次隔离运行除实际耗时外 48 个正式输出文件逐字节一致，硬约束最大残差不超过 $10^{-6}$。所得结论属于有限搜索下的 scenario-feasible 结果，不声称全局最优。
\end{abstract}

\noindent\textbf{关键词：}异构 AI 任务；算力调度；Huber 回归；储能协同；碳感知调度；情景分析
